%&bericht

%%%%%%%%%%%%%%%%%%%%%%%%%%%%%%%%%%%%%%%%%%%%%%%%%%%%%%%%%%%%%%%%%%%%%%%%%%%%%%%
%% Descr:       Vorlage für Berichte der DHBW-Karlsruhe
%% Author:      Prof. Dr. Jürgen Vollmer, juergen.vollmer@dhbw-karlsruhe.de
%% $Id: bericht.tex,v 1.25 2020/03/13 15:07:45 vollmer Exp $
%%  -*- coding: utf-8 -*-
%%%%%%%%%%%%%%%%%%%%%%%%%%%%%%%%%%%%%%%%%%%%%%%%%%%%%%%%%%%%%%%%%%%%%%%%%%%%%%%

\documentclass[
   ngerman          % neue deutsche Rechtschreibung
  ,a4paper          % Papiergrösse
% ,twoside          % Zweiseitiger Druck (rechts/links)
% ,10pt             % Schriftgrösse
  ,11pt
% ,12pt
  ,pdftex
%  ,disable         % Todo-Markierungen auschalten
]{report}

% Bitte die Codierung Ihrer Dateien auswählen:
% \usepackage[latin1]{inputenc}    % Für UNIX mit ISO-LATIN-codierten Dateien
% \usepackage[applemac]{inputenc}  % Für Apple Mac
% \usepackage[ansinew]{inputenc}   % Für Microsoft Windows
\usepackage[utf8]{inputenc}        % UTF-8 codierte Dateien
                                   % Dieses Dokument ist unter Unix erstellt, daher
                                   % wird diese Input-Codierung benutzt.

\usepackage{bericht}
\usepackage{amsmath}
%% ACHTUNG, wenn man eine eigene Formatdatei (bericht.fmt) benutzt, werden Änderungen an bericht.sty
%% erst wirksam, wenn die Format-Datei neu erzeugt wurde!!!
%% Genauer alle Änderungen, die textuell vor der nächsten Zeile ".... endofdump...." stehen
%% werden erst wirksam, wenn die Formatdatei neu erzeugt wurde
\csname endofdump\endcsname

%%%%%%%%%%%%%%%%%%%%%%%%%%%%%%%%%%%%%%%%%%%%%%%%%%%%%%%%%%%%%%%%%%%%%%%%%%%%%%%
%% Angaben zur Arbeit
%%%%%%%%%%%%%%%%%%%%%%%%%%%%%%%%%%%%%%%%%%%%%%%%%%%%%%%%%%%%%%%%%%%%%%%%%%%%%%%

\newcommand{\Autor}{Max Mustermann}
\newcommand{\MatrikelNummer}{4711}
\newcommand{\Kursbezeichnung}{tinf17b3}

\newcommand{\FirmenName}{Firmenname}
\newcommand{\FirmenStadt}{Stadt}
%SELF \newcommand{\FirmenLogoDeckblatt}{\fbox{\includegraphics[width=3cm]{lion}}}

% Falls es kein Firmenlogo gibt:
%  \newcommand{\FirmenLogoDeckblatt}{}

\newcommand{\BetreuerFirma}{Titel Vorname Nachname}
\newcommand{\BetreuerDHBW}{Titel Vorname Nachname}

%%%%%%%%%%%%%%%%%%%%%%%%%%%%%%%%%%%%%%%%%%%%%%%%%%%%%%%%%%%%%%%%%%%%%%%%%%%%%%%%%%%%%

% Wird auf dem Deckblatt und in der Erklärung benutzt:
\newcommand{\Was}{Projekt-/Studien-/Bachleorarbeit}
%\newcommand{\Was}{Projektrarbeit}
%\newcommand{\Was}{Studienarbeit}
%\newcommand{\Was}{Bachleorarbeit}

%%%%%%%%%%%%%%%%%%%%%%%%%%%%%%%%%%%%%%%%%%%%%%%%%%%%%%%%%%%%%%%%%%%%%%%%%%%%%%%%%%%%%

\newcommand{\Titel}{\LaTeX-Vorlage für diverse Ausarbeitungen\\--\\oder so ähnlich}
\newcommand{\AbgabeDatum}{1. April 2090}

\newcommand{\Dauer}{12 Wochen}

% \newcommand{\Abschluss}{Bachelor of Engineering}
\newcommand{\Abschluss}{Bachelor of Science}

\newcommand{\Studiengang}{Informatik / Informationstechnik}
% \newcommand{\Studiengang}{Informatik / Angewandte Informatik}

\hypersetup{%%
  pdfauthor={\Autor},
  pdftitle={\Titel},
  pdfsubject={\Was}
}

%%%%%%%%%%%%%%%%%%%%%%%%%%%%%%%%%%%%%%%%%%%%%%%%%%%%%%%%%%%%%%%%%%%%%%%%%%%%%%%

% Wenn \includeonly{..} benutzt wird, werden nur diese Kaptitel ausgegeben.
\includeonly{
  abk
 ,kapitel1
 ,kapitel2
 ,changelog
}

%%%%%%%%%%%%%%%%%%%%%%%%%%%%%%%%%%%%%%%%%%%%%%%%%%%%%%%%%%%%%%%%%%%%%%%%%%%%%%%

% Benutzt man das "biblatex"-Paket, dann muß das hier stehen:
% siehe auch die mit BIBLATEX markierten Zeilen in bericht.sty
\bibliography{bericht}

\begin{document}

%%%%%%%%%%%%%%%%%%%%%%%%%%%%%%%%%%%%%%%%%%%%%%%%%%%%%%%%%%%%%%%%%%%%%%%%%%%%%%%

\begin{titlepage}
\begin{center}
\vspace*{-2cm}
%SELF \FirmenLogoDeckblatt\hfill\includegraphics[width=4cm]{dhbw-logo}\\[2cm]
{\Huge \Titel}\\[1cm]
{\Huge\scshape \Was}\\[1cm]
{\large für die Prüfung zum}\\[0.5cm]
{\Large \Abschluss}\\[0.5cm]
{\large des Studienganges \Studiengang}\\[0.5cm]
{\large an der}\\[0.5cm]
{\large Dualen Hochschule Baden-Württemberg Karlsruhe}\\[0.5cm]
{\large von}\\[0.5cm]
{\large\bfseries \Autor}\\[1cm]
{\large Abgabedatum \AbgabeDatum}
\vfill
\end{center}
\begin{tabular}{l@{\hspace{2cm}}l}
Bearbeitungszeitraum	         & \Dauer 			\\
Matrikelnummer	                 & \MatrikelNummer		\\
Kurs			         & \Kursbezeichnung		\\
Ausbildungsfirma	         & \FirmenName			\\
			         & \FirmenStadt			\\
Betreuer der Ausbildungsfirma	 & \BetreuerFirma		\\
Gutachter der Studienakademie	 & \BetreuerDHBW		\\
\end{tabular}
\end{titlepage}

%%%%%%%%%%%%%%%%%%%%%%%%%%%%%%%%%%%%%%%%%%%%%%%%%%%%%%%%%%%%%%%%%%%%%%%%%%%%%%%

%%%%%%%%%%%%%%%%%%%%%%%%%%%%%%%%%%%%%%%%%%%%%%%%%%%%%%%%%%%%%%%%%%%%%%%%%%%%%%%
%% Descr:       Vorlage für Berichte der DHBW-Karlsruhe, Erklärung
%% Author:      Prof. Dr. Jürgen Vollmer, vollmer@dhbw-karlsruhe.de
%% $Id: erklaerung.tex,v 1.11 2020/03/13 14:24:42 vollmer Exp $
%% -*- coding: utf-8 -*-
%%%%%%%%%%%%%%%%%%%%%%%%%%%%%%%%%%%%%%%%%%%%%%%%%%%%%%%%%%%%%%%%%%%%%%%%%%%%%%%

% In Bachelorarbeiten muss eine schriftliche Erklärung abgegeben werden.
% Hierin bestätigen die Studierenden, dass die Bachelorarbeit, etc.
% selbständig verfasst und sämtliche Quellen und Hilfsmittel angegeben sind. Diese Erklärung
% bildet das zweite Blatt der Arbeit. Der Text dieser Erklärung muss auf einer separaten Seite
% wie unten angegeben lauten.

\newpage
\thispagestyle{empty}
\begin{framed}
\begin{center}
\Large\bfseries Erklärung
\end{center}
\medskip
\noindent
% siehe §5(3) der \enquote{Studien- und Prüfungsordnung DHBW Technik} vom 29.\,9.\,2017 und Anhang 1.1.13
Ich versichere hiermit, dass ich meine \Was mit dem Thema:
\enquote{\Titel}
selbstständig verfasst und keine anderen als die angegebenen Quellen und Hilfsmittel benutzt habe. Ich versichere zudem, dass die eingereichte elektronische Fassung mit der gedruckten Fassung übereinstimmt.
\vspace{3cm}
\noindent
\underline{\hspace{4cm}}\hfill\underline{\hspace{6cm}}\\
Ort~~~~~Datum\hfill Unterschrift\hspace{4cm}
\end{framed}

\vfill
\emph{Sofern  vom Dualen Partner ein Sperrvermerk gewünscht wird, ist folgende Formulierung
zu verwenden:}
\begin{framed}
\begin{center}
\Large\bfseries Sperrvermerk
\end{center}
\medskip
\noindent
Der Inhalt dieser Arbeit darf weder als Ganzes noch in Auszügen Personen
außerhalb des Prüfungsprozesses und des Evaluationsverfahrens zugänglich gemacht
werden, sofern keine anderslautende Genehmigung vom Dualen Partner vorliegt.
\end{framed}

%%%%%%%%%%%%%%%%%%%%%%%%%%%%%%%%%%%%%%%%%%%%%%%%%%%%%%%%%%%%%%%%%%%%%%%%%%%%%%%
\endinput
%%%%%%%%%%%%%%%%%%%%%%%%%%%%%%%%%%%%%%%%%%%%%%%%%%%%%%%%%%%%%%%%%%%%%%%%%%%%%%%

%%%%%%%%%%%%%%%%%%%%%%%%%%%%%%%%%%%%%%%%%%%%%%%%%%%%%%%%%%%%%%%%%%%%%%%%%%%%%%%

\begin{abstract}
Dieses \LaTeX-Dokument kann als Vorlage für einen Praxis- oder Projektbericht, eine Studien- oder
Bachelorarbeit dienen.

Zusammengestellt von Prof.\,Dr.\,Jürgen Vollmer \email{juergen.vollmer@dhbw-karlsruhe.de}\\
\url{https://www.karlsruhe.dhbw.de}. Die jeweils aktuellste Version dieses \LaTeX-Paketes ist immer
auf der \emph{FAQ-Seite} des Studiengangs Informatik zu finden:
\url{https://www.karlsruhe.dhbw.de/inf/studienverlauf-organisatorisches.html} $\to$ \emph{Formulare und Vorlagen}.

\centering Stand \verb+$Date: 2020/03/13 15:07:45 $+
\end{abstract}

\newpage
\tableofcontents           % Inhaltsverzeichnis hier ausgeben
\listoffigures             % Liste der Abbildungen
\listoftables              % Liste der Tabellen
\lstlistoflistings         % Liste der Listings
\listofequations           % Liste der Formeln

% Jetzt kommt der "eigentliche" Text
%%%%%%%%%%%%%%%%%%%%%%%%%%%%%%%%%%%%%%%%%%%%%%%%%%%%%%%%%%%%%%%%%%%%%%%%%%%%%%
%% Descr:       Vorlage für Berichte der DHBW-Karlsruhe, Datei mit Abkürzungen
%% Author:      Prof. Dr. Jürgen Vollmer, vollmer@dhbw-karlsruhe.de
%% $Id: abk.tex,v 1.4 2017/10/06 14:02:03 vollmer Exp $
%% -*- coding: utf-8 -*-
%%%%%%%%%%%%%%%%%%%%%%%%%%%%%%%%%%%%%%%%%%%%%%%%%%%%%%%%%%%%%%%%%%%%%%%%%%%%%%%

\chapter*{Abkürzungsverzeichnis}                   % chapter*{..} -->   keine Nummer, kein "Kapitel"
						         % Nicht ins Inhaltsverzeichnis
% \addcontentsline{toc}{chapter}{Akürzungsverzeichnis}   % Damit das doch ins Inhaltsverzeichnis kommt

% Hier werden die Abkürzungen definiert
\begin{acronym}[DHBW]
  % \acro{Name}{Darstellung der Abkürzung}{Langform der Abkürzung}
 \acro{Abk}[Abk.]{Abkürzung}

 % Folgendes benutzen, wenn der Plural einer Abk. benöigt wird
 % \newacroplural{Name}{Darstellung der Abkürzung}{Langform der Abkürzung}
 \newacroplural{Abk}[Abk-en]{Abkürzungen}


 \acro{RSA}[RSA]{Rivest-Shamir-Adleman}
 \acro{AES}[AES]{Advanced Encryption Standard}

\end{acronym}
              % Abkürzungsverzeichnis

\chapter{Grundlagen}

In diesem Kapitel werden die theoretischen Grundlagen der Kryptologie, Mathematik, Zahlentheorie und der IT-Sicherheit erläutert, die in dieser Arbeit eine Rolle spielen. Aus diesen sollen die grundlegenden Funktionen eines kleptographischen Angriffes und dessen Folgen abgeleitet werden.

\section{Kryptologie}
    Die Kryptologie ist die wissenschaftliche Disziplin für den Schutz von Daten. Unter ihr stehen die zwei Felder der Kryptographie und der Kryptoanalyse.

    \subsection{Kryptographie}
        Die Kryptologie befasst sich mit der Entwicklung von Verfahren und Techniken für den sicheren Austausch von Daten. Dabei stehen zwei Eigenschaften \cite{BSW.2015} im Fokus:

        \subsubsection{Eigenschaften}
            \paragraph{Geheimhaltung}
                Durch Geheimhaltung sollen, bei der Übertragung von Daten zwischen Teilnehmern, Unbeteiligte keine Erkenntnisse über den Inhalt erlangen. Dies kann durch physikalische oder organisatorische Maßnahmen erreicht werden, wobei Unbeteiligten der Zugang zu den übertragenen Daten verwehrt wird. Diese Maßnahmen sind sinnvoll bei der Übergabe der Daten in einer nicht digitalen Welt. Bei der Kommunikation in digitalen Netzen, wie u.a. dem Internet, sind diese Maßnahmen nur schwer zu implementieren. Dies gilt nicht für kryptographische Maßnahmen. Dabei ist es nicht mehr das Ziel, Unbeteiligten den Zugang zu den übertragenen Daten zu erschweren, sondern den Inhalt der Daten während der Übertragung zu verschlüsseln. Dadurch soll es Unbeteiligten nahezu unmöglich sein, aus den mitgehörten oder abgefangenen Daten, Rückschlüsse auf deren Inhalt zu erlangen.
                
            \paragraph{Authentifikation}
                Durch Authentifikation soll es den Teilnehmern einer Kommunikation möglich sein, die anderen Teilnehmer und empfangene Nachrichten zweifelsfrei identifizieren und zuweisen zu können. Hierbei spielen Signaturverfahren eine wichtige Rolle, da kein Geheimnis benötigt wird um einen Teilnehmer zu authentifizieren. Dabei können sich Teilnehmer durch das Wissen oder den Besitz eines Geheimnisses (Passwort, Zertifikat, Schlüssel) authentifizieren. \cite{BSW.2015}

        Nur wenn beide Eigenschaften gegeben sind ist eine Übertragung von Daten als sicher anzusehen. Falls die Geheimhaltung fehlt, kann der Inhalt durch Sniffing mitgelesen werden. Falls die Authentifikation der Teilnehmer fehlt, können sich Unbeteiligte als "echte" Teilnehmer ausgeben und somit die Daten an ihrem Endpunkt entschlüsseln.

        \subsubsection{Verschlüsselungsverfahren}
        % Was sind Verschlüsselungsverfahren, Ziel, Historie, momentaner Stand, State of the Art (AES, RSA )
            \paragraph{Asymmetrische Verschlüsselung}
            % Hier angesprochene Verfahren, unterschiede, Definition / Abgrenzung zu symmetrischen Verfahren

            \paragraph{Hybride Verfahren}
                % Meiste Verfahren heute hybride Verfahren. Asymmetrisch zum Schlüsselaustausch. Symmetrisch zum Verschlüsseln. Laufzeit von Asymmetrischen Verfahren. Wenn Schlüsselaustausch sicher Geheimhaltung verletzt.

        \subsubsection{Angriffe}
            % Bedeutung von Brute-Force bei modernen kryptographischen Verfahren

            % Ziel der Angriffe / Bedeutung von Schlüsseln
            
            \paragraph{Known Cipher Attack}

            \paragraph{Known Plaintext Attack}

            \paragraph{Chosen Plaintext Attack}

            \paragraph{Chosen Cipher Attack}


\section{Mathematik}
    Mathematische Probleme stellen die Grundlage für moderne Kryptographie. 

    \subsection{Diskreter Logarithmus}
        Bei der Bestimmung des Logarithmus wird der Exponent (hier: $x$) gesucht, welcher mit einer bekannten Zahl als Basis $z$, eine weitere bekannte Zahl $y$ ergibt.
        \begin{equation}
            z^{x} = y
            \eqlabel{eq-log}{Normaler Logarithmus}
        \end{equation}

        Der diskrete Logarithmus bezieht hier auf die Berechnung des Logarithmus in ein Gruppe. Diese Gruppe bildet sich aus der Rechnung mit Restklassen (modulo). Dadurch entsteht folgendes Problem, bei der die Variable $x$ gesucht ist und alle anderen Variablen bekannt sind.
        \begin{equation}
            z^{x} \pmod n \equiv y
            \eqlabel{eq-dis-log}{Diskreter Logarithmus}
        \end{equation}
        
        Hierbei ist in der Notation zu beachten, dass sich durch das Rechnen auf mit einer Gruppe, Äquivalenzklassen ($\equiv$) bilden. Diese entsprechen den Restklassen des Rechnen mit Modulo. $n$ ist die Mächtigkeit der Aquivalenzklassen.

        Die Bestimmung von $x$ in \ref{eq-dis-log} wird als Problem des diskreten Logarithmus bezeichnet. Mit der Komplexität wird sich in den Grundlagen der Komplexitätstheorie beschäftigt.

        Dabei ist die Umkehrfunktion, des diskreten Logarithmus $f(x)$ \ref{eq-dis-log}, mathematisch einfach zu berechnen. Diese Umkehrfunktion entspricht der diskreten Exponentialfunktion:
        \begin{equation}
            f^{-1}(x) = z^{x} \pmod n \equiv y
            \eqlabel{eq-dis-exp}{Diskretere Exponentialfunktion}
        \end{equation}
        Hierbei sind $z,x,n$ gegeben und $y$ gesucht.

    \subsection{Faktorisierung}


    \subsection{Effiziente Berechnung der diskreten Exponentialfunktion}
        In der Kryptographie werden große Zahlen genutzt, um die Sicherheit der verwendeten Algorithmen zu gewährleisten. Hierfür wird als Beispiel angenommen, dass als Basis $z$ eine 256-bit Lange Zahl hoch einem 300-bit langem Exponenten $x$ genommen werden soll. Hierbei ist $n$ 1024-bit lang. 

        Wenn mann nun $z$ in Byte berechnet wäre dies eine 32 Byte lange Zahl.

        $x$ entspricht einer ungefähr 90. stelligen Zahl. 

        \begin{equation}
            z^{10^{90}} \pmod n \equiv y
            \eqlabel{eq-dis-exp-lN}{Diskretere Exponentialfunktion mit großen Zahlen}
        \end{equation}

        Eine numerische Berechnung von $ z^{10^{90}} $ ist aufgrund von begrenzten Ressourcen nicht möglich. 

        Jedoch kann man sich die diskrete Eigenschaft dieser Problems sich zu nutzte machen. Hierfür können Verfahren, wie Square-and-Multiply zusammen mit der Restklassenberechnung genutzt werden. Dadurch lassen sich auch großzahlige Exponenten berechnen. Hierfür soll ein einfaches Beispiel gegeben werden:
        \begin{equation}
            37^{52} \pmod {128} \equiv y
            \eqlabel{eq-dis-exp-bspOne}{Diskretere Exponentialfunktion mit großen Zahlen Beispiel-Eins}
        \end{equation}
        Bei Betrachtung der Äquivalenzgleichung fällt auf, dass $37^{52}$ eine große Zahl ergibt. Jedoch wird diese Zahl noch $ x \pmod 128$ gerechnet. Dadurch liegt das Ergebnis in einem Zahlenraum von:
        \begin{equation}
            y \in \mathbb{N} \mid 0 \le y < 128 
            \eqlabel{eq-dis-exp-numSpace}{Diskretere Exponentialfunktion in Zahlenraum}
        \end{equation}

        Auf Grundlage der Potenzgesetze wird $37^{52}$ nun zerlegt. 
        \begin{equation}
        \begin{aligned}
            52 = 32 + 16 + 4 = 2^{5} + 2^{4} + 2^{2} \\
            37^{52} \pmod {128} \equiv 37^{2^{5}} * 37^{2^{4}} * 37^{2^{2}} \\
            \equiv 37^{2^{5}} \pmod {128} * 37^{2^{4}} \pmod {128} * 37^{2^{2}} \pmod {128}
        \end{aligned}
        \eqlabel{eq-dis-exp-bspTwo}{Diskretere Exponentialfunktion mit großen Zahlen Beispiel-Zwei}
        \end{equation}

        Die einzelnen Bestandteile werden dann iterativ berechnet und durch Multiplikation zusammengefasst (siehe \ref{eq-dis-exp-bspTwo}). Dies wird als Square-and-Multiply-Verfahren bezeichnet.
        \begin{equation}
        \begin{aligned}
            37^{2^{2}} \pmod {128} \equiv (37^{2^{1}} \pmod {128})^{2} \\
            37^{2^{3}} \pmod {128} \equiv (37^{2^{2}} \pmod {128})^{2} \\
            37^{2^{4}} \pmod {128} \equiv (37^{2^{3}} \pmod {128})^{2} \\
            37^{2^{5}} \pmod {128} \equiv (37^{2^{4}} \pmod {128})^{2}
        \end{aligned}
        \eqlabel{eq-dis-exp-bspThree}{Diskretere Exponentialfunktion mit großen Zahlen Beispiel-Drei}
        \end{equation}
        
        \subsubsection{Allgemein}
        Gegeben mit gesucht $y$:
        \begin{equation}
            z^{x} \pmod n \equiv y
        \end{equation}

        Zerlegung von $x$ eine Summe von Zweierpotenzen:
        \begin{equation}
            x = 2^{0} + 2^{1} + 2^{2} + ...
        \end{equation}

        Dabei bilden die binären Logarithmen der einzelnen Zweierpotenzen die Menge $\mathbb{K}$.

        Berechnung der einzelnen Faktoren durch iteratives Square-and-Multiply-Verfahren. Dies wird bis $f(max(\mathbb{K}))$ berechnet. $max(\mathbb{K})$ steht hier für das Element von $\mathbb{K}$, mit dem größten Wert.
        \begin{equation}
            f(i+1) = f(i)^{2} \pmod n \mid f(1) = z^{1} \pmod n 
        \end{equation}

        Zuletzt wird das Produkt, aller Ergebnisse von $f(x)$ für die Elemente der Menge $\mathbb{K}$, gebildet. Dabei gilt:

        \begin{equation}
            \prod_{k \in \mathbb{K}} f(k) \equiv z^{x} \pmod n \equiv y
        \end{equation}





\chapter{RSA}
\ac{RSA} ist ein kryptographisches Verfahren, welches zu den Public-Key-Verfahren gehört. Der Verfahren wurde von R. Rivest, A. Shamir und L. Adleman entwickelt und trägt deshalb ein Anagram der Erfinder als Namen.

\section{Ablauf} \label{rsa-ablauf}
    \paragraph{Ausgangszenario} Teilnehmer A will über ein öffentliches Netz sicher mit anderen Teilnehmer kommunizieren. 

    \paragraph{Schlüsselgeneration} Damit andere Teilnehmer geheime Nachrichten schicken können muss A sich ein Schlüsselpaar generieren. Dafür wählt er zwei zufällige und große Primzahlen: $P$ und $Q$. 

    Das Produkt von $P$ und $Q$ bildet $N$, welche den Modulo / den Zahlenraum für weitere mathematische Operationen festlegt: 
        \begin{equation}
            N = P \cdot Q
            \eqlabel{eq-rsa-npq}{RSA - Primfaktoren}
        \end{equation}

    Daraufhin berechnet der Teilnehmer die Eulersche $\phi$-Funktion von $P$ und $Q$:
        \begin{equation}
            \phi(N) = (P-1) \cdot (Q-1)
            \eqlabel{eq-rsa-phi_n}{RSA - Eulersche $\phi$-Funktion}
        \end{equation})
    
    Der öffentliche Schlüssel $E$ ist dann eine zu $\phi(N)$ teilerfremde Zahl. Man kann dies auch vereinfachen und eine Fermat'sche Primzahl für $E$ verwenden:
        \begin{equation}
            2^{2^{N}}+1 \mid N \in \{0,1,2,3,4\}
            \eqlabel{eq-rsa-ferma}{Fermat'sche Primzahl}
        \end{equation}
    
    Der größte gemeinsame Teiler von $E$ und $\phi(n)$ ergibt $1$, wobei sich der private Schlüssel $D$ aus der Vielfachsummendarstellung ergibt. Die Vielfachsummendarstellung wird dem euklidischen Algorithmus entnommen.  Dabei gilt:
        \begin{equation}
            c \cdot \phi(N) + E \cdot D \equiv 1
            \eqlabel{eq-rsa-calc_d}{RSA - Berechnung des privaten Schlüssels}
        \end{equation}
    
    Danach besitzt der Teilnehmer A ein öffentlichen Schlüssel $E$ und einen privaten Schlüssel $D$. Er kann nun den öffentlichen Schlüssel $E$ zusammen mit $N$ veröffentlichen. $E$ zusammen mit $N$ und $D$ sind ein Schlüsselpaar welches in einen öffentlichen Teil und ein privaten Teil geteilt wird. Dabei herrscht eine eindeutige Zuordnung zwischen den Teilen.
    
    \subsection{Verschlüsselung} \label{sec-RSA-crypt}
        Wenn Teilnehmer B eine verschlüsselte Nachricht $M$ an Teilnehmer A senden will, braucht er hierfür den öffentlichen Schlüssel von A. $M$ muss aber dabei kleiner als $N$ sein. Dabei B  verschlüsselt wie folgt, wobei $C$ der resultierende Geheimtext ist.

        \begin{equation}
            M^{E_{A}} \pmod N = C
            \eqlabel{eq-rsa-encrypt}{RSA - Verschlüsselung einer Nachricht}
        \end{equation}

        Wenn Teilnehmer A den Geheimtext $C$ von B erhält, entschlüsselt er diesen mit seinem privaten Schlüssel. Dadurch berechnet er die von Teilnehmer B verschlüsselte Nachricht $M$.

        \begin{equation}
            C^{D_{A}} \pmod N = M
            \eqlabel{eq-rsa-decrypt}{RSA - Entschlüsselung eines Geheimtextes}
        \end{equation}

        Da $D$ privat und eindeutig für den öffentlichen Teil des Schlüsselpaars ist, können nur Teilnehmer, die über $D$ verfügen, Nachrichten entschlüsseln, die mit dem zugehörigen öffentlichen Teil verschlüsselt wurden.
    
    \subsection{Signatur} \label{sec-RSA-sign}
        Das RSA-Verfahren ermöglicht auch das Signieren von Nachrichten. Dies wird z.B. im Internet genutzt um Teilnehmer zu authentifizieren und Nachrichtenintegrität zu sichern. RSA wird jedoch meist nur auf kleinere Nachrichten wie Fingerabdrücke angewandt.

        Signatur $S$ einer Nachricht $M$ von Teilnehmer A:
        \begin{equation}
            M^{D_{A}} \pmod N = S
            \eqlabel{eq-rsa-sign}{RSA - Signieren einer Nachricht}
        \end{equation}

        Zum Überprüfen der Echtheit braucht der Teilnehmer B den öffentlichen Schlüssel von A:
        \begin{equation}
            S^{E_{A}} \pmod N = M
            \eqlabel{eq-rsa-check_sign}{RSA - Prüfen einer Signatur}
        \end{equation}

        Eine Signatur kann eindeutig Teilnehmer A zugeordnet werden, da nur ein Teilnehmer im Besitz von $d_{A}$ eine Signatur einer Nachricht erstellen kann, die mit dem öffentlichen Teil des Schlüsselpaars von A überprüft werden kann.
    
\section{Annahmen} \label{sec-RSA-assertion}
        Das Ergebnis einer \ac{RSA}-Verschlüsselung wird als Zufallszahl betrachtet. Grundlage hierfür ist, dass RSA mit sicheren Schlüsseln und Implementation, einer Einweg-Funktion gleicht, wenn der private Schlüssel nicht bekannt ist. Das Ergebnis einer \ac{RSA}-Verschlüsselung ist folgendermaßen die Generierung einer Zufallszahl $C$:
        \begin{equation}
            C \in [1,N)
            \eqlabel{eq-rsa-random}{RSA-Verschlüsselung EInwegfunktion}
        \end{equation}
        % Kann C 0 oder 1 sein? 
        Diese Annahme ist für die mathematische Betrachtung von möglichen Grenzfälle später relevant. Im kryptographischen Kontext, kann natürlich die Einwegfunktion mit Kenntnis des privaten Schlüssels effizient umgekehrt werden.

\section{Sicherheit}
    Die Sicherheit des RSA-Verfahrens basiert auf zwei mathematischen Problemen, welche unter Aufwand endlicher Ressourcen, nicht gelöst werden können. Hierbei wird sich sowohl auf RSA-gestützte Verschlüsselungs- und Signaturverfahren bezogen.
    Diese zwei Probleme sind:
    \begin{itemize}
        \item Faktorisierung einer bekannten Zahl, welche das Produkt zweier großer Primzahlen ist. Im Kontext von RSA ist diese Zahl mit $N$ repräsentiert.
        \item Bestimmung des diskreten Logarithmus. Bei RSA wäre dies die Bestimmung von 
        \begin{equation}
            D \mid M^{D} \equiv C \pmod N .
        \end{equation}
    \end{itemize}

    Für die Sicherheit der Public-Key-Verschlüsselung von RSA, spielt nicht effiziente Berechenbarkeit der Faktorisierung die Hauptrolle. Falls mit RSA signiert werden soll, ist zusätzlich die Unberechenbarkeit des diskreten Logarithmus wichtig. Ansonsten könnte der private Schlüssel bestimmt werden.

    \paragraph{Angriffe auf RSA}
    % Momentane Angriffe auf RSA
    % RSA ist momentan als sicher bekannt

    % Fermatsche Faktorisierung
    % Gemeinsame Primzahlen


% Ab hier beginnt der Anhang
\appendix
\addcontentsline{toc}{chapter}{Anhang}

\addcontentsline{toc}{chapter}{Index}
\printindex

\addcontentsline{toc}{chapter}{Literaturverzeichnis}




% Haben Sie das "biblatex"-Paket nicht installiert, benutzen Sie folgendes:
% Ohne das "biblatex"-Paket (s. bericht.sty) produziert folgendes
% "deutsche" Zitate in Literaturverzeichnissen gemaß der Norm DIN 1505,
% Teil 2 vom Jan. 1984.
% Die Zitatmarken werden alphabetisch nach Verfassern
% sortiert und sind durch abgekürzte Verfasserbuchstaben plus
% Erscheinungsjahr in eckigen Klammern gekennzeichnet.

% \bibliographystyle{alpha}
% \bibliography{bericht}

%%%%%%%%%%%%%%%%%%%%%%%%%%%%%%%%%%%%%%%5
% BIBLATEX
% Benutzt man das "biblatex"-Paket, muß man folgendes schreiben:
\def\refname{Literaturverzeichnis}
\printbibliography
%%%%%%%%%%%%%%%%%%%%%%%%%%%%%%%%%%%%%%%5


\end{document}
